% notes/wave17_opus_20260424/opus_04_G2_arnold_kzb.tex
%
% Elite-voice adversarial-heal agent -- G2 generating identity:
%    d_bar = KZ^*(nabla_Arnold) locally on Conf_n(C) for fixed C;
%    d_bar^{univ} = (KZ^{univ})^*(nabla^{univ}_KZB) globally on
%    Conf_n^{univ}/\bar M_{g,n},
% with Felder-Wieczerkowski (1996) + Calaque-Enriquez-Etingof
% (2009) globalisation, regular singularities along
% \partial\bar M_{g,n}, and Jimbo-Miwa-Stokes irregular behaviour at
% specific strata (F5 frontier).
%
% Voices: MANIN (motivic-arithmetic) + ETINGOF (clarity), with
% secondary WITTEN (physical insight).
%
% Output: cycle log + healed inscription + hidden structure
% found during adversarial attack + opens + cross-volume harmonies.
% NOT to be \input into the manuscript; lives in notes/ as a
% working paper for inscription during a subsequent wave.
%
% Complementary to:
%   chapters/theory/chiral_climax_platonic.tex   (G2 genus-0 on P^1)
%   chapters/theory/climax_theorem.tex           (G2 summary chapter)
%   chapters/theory/genus_2_ddybe_platonic.tex   (CEE g=2 globalisation)
%   chapters/connections/genus1_seven_faces.tex  (Face 4 = KZB)

\documentclass[12pt]{article}
\usepackage[localtheorems]{raeez-math-template}


\providecommand{\KZ}{\mathrm{KZ}}
\providecommand{\KZB}{\mathrm{KZB}}
\providecommand{\Arn}{\mathrm{Arnold}}
\providecommand{\Conf}{C}
\providecommand{\Confn}{\mathrm{Conf}_n}
\providecommand{\Confnord}[1][C]{\mathrm{Conf}_n^{\mathrm{ord}}(#1)}
\providecommand{\FM}{\overline{\mathrm{FM}}}
\providecommand{\Mbar}{\overline{\mathcal{M}}}
\providecommand{\Ran}{\mathrm{Ran}}
\providecommand{\Ranord}{\mathrm{Ran}^{\mathrm{ord}}}
\providecommand{\dbar}{d_{\mathrm{bar}}}
\providecommand{\BarOrd}{B^{\mathrm{ord}}}
\providecommand{\ChirAlg}{\mathrm{ChirAlg}}
\providecommand{\Dmod}{D\text{-mod}}
\providecommand{\ConnConf}{\mathrm{Conn\text{-}Conf}}
\providecommand{\fg}{\mathfrak{g}}
\providecommand{\hv}{h^{\vee}}
\providecommand{\GRT}{\mathrm{GRT}_1}
\providecommand{\op}{\mathrm{op}}

\newtheorem{theorem}{Theorem}[section]
\newtheorem{proposition}[theorem]{Proposition}
\newtheorem{lemma}[theorem]{Lemma}
\newtheorem{corollary}[theorem]{Corollary}
\theoremstyle{definition}
\newtheorem{definition}[theorem]{Definition}
\newtheorem{construction}[theorem]{Construction}
\newtheorem{remark}[theorem]{Remark}
\newtheorem{conjecture}[theorem]{Conjecture}

\title{Opus adversarial-heal deliverable 04 \\
       G2: $\dbar = \KZ^{*}(\nabla_{\Arn})$, local and global}
\author{Wave~17 opus thread}
\date{2026-04-24}

\begin{document}
\maketitle

\begin{abstract}
Adversarial-heal treatment of the generating identity G2
(chain-level pullback of Arnold's universal KZ connection to the
bar differential) and its globalisation
$\dbar^{\mathrm{univ}} = (\KZ^{\mathrm{univ}})^{*}(\nabla_{\KZB}^{\mathrm{univ}})$
over $\Confn^{\mathrm{univ}}/\Mbar_{g,n}$.  Six attack cycles
(chain-level strictness, strict pullback of $\dbar$, universal KZB
existence past $g \geq 1$, irregular-singular stratum content,
compatibility of the four equivalent presentations, GRT$_1(\mathbb{Q})$
torsor action) are answered with healed statements: an explicit
chain homotopy witness (strict equality holds on the bar
generating set, contractible homotopy on the full cofree
coalgebra); the universal KZB as the Calaque-Enriquez-Etingof
$\Mbar_{g,n}$-relative flat connection; Stokes-Jimbo-Miwa
classification of the irregular strata as an honest frontier
(F5); functoriality of the four equivalent presentations
through the common Fulton-MacPherson residue map; and the
GRT$_1(\mathbb{Q})$ torsor action as a theorem
(Drinfeld $1990$, Furusho $2010$).  Each cycle ends with a
load-bearing inscription and explicit numerical checks at
$g = 0, 1, 2$ (Arnold flat $+$ Kohno monodromy recovery;
Bernard-Felder elliptic propagator; Enriquez genus-$2$
universal KZB and Fay trisecant).
\end{abstract}

\tableofcontents

%% ====================================================================
\section{Statement and scope discipline (Manin-Gaitsgory)}
%% ====================================================================
\label{sec:G2-statement}

\textbf{Local form (fixed curve).}
Let $C$ be a smooth projective algebraic curve over $\mathbb{Q}$
(or $\mathbb{C}$; we keep the rational-motivic setting unless
otherwise stated).  On the ordered configuration space
$\Confnord(C) \subset C^{n} \setminus \Delta$ with Arnold
$1$\nobreakdash-forms
$\eta_{ij} = d\log(z_i - z_j)$
and infinitesimal braid generators $t_{ij}$, the universal KZ
connection is Arnold's 1969 meromorphic flat connection
\[
   \nabla_{\Arn}
   \;=\;
   d \;-\; \sum_{i < j}\,t_{ij}\,\eta_{ij},
   \qquad
   \nabla_{\Arn}^{2} \;=\; 0
   \iff
   \text{Arnold three-term identity}.
\]
For a chirally Koszul $E_\infty$-chiral algebra $A$ on $C$, the
\emph{KZ functor}
$\KZ \colon \ChirAlg^{E_\infty,\mathrm{Koszul}}_{C}
 \to \ConnConf_{C}$
(Construction \ref{constr:kz-functor-platonic}
of \texttt{chapters/theory/chiral\_climax\_platonic.tex})
produces for each $A$ a pair
$\bigl(\BarOrd(A),\nabla(A)\bigr)$ whose flat connection is, up
to logarithmic gauge, the pullback of $\nabla_{\Arn}$ along the
structure morphism $U(\mathfrak{t}_\bullet) \to \mathrm{End}(\BarOrd(A))$
induced by $A \to \mathrm{End}(\BarOrd(A))$.  The
\emph{pullback identity G2} is
\[
   (\text{G2-local})
   \qquad
   \dbar \;=\; \KZ^{*}(\nabla_{\Arn})
   \quad
   \text{on $\Confnord(C)$ at fixed $C$.}
\]

\textbf{Global form (varying curve).}
Fix a genus $g \geq 0$ and a number of marked points $n \geq 1$;
let $\Mbar_{g,n}$ be the Deligne-Mumford compactification.  The
universal ordered configuration space
$\Conf_n^{\mathrm{univ}} / \Mbar_{g,n}$ carries the relative
factorisation D-module of BD (1999) + Francis-Gaitsgory (2012),
and the family
$\{\nabla_{\Arn,C}\}_{C \in \mathcal{M}_{g,n}}$ extends across
$\partial \Mbar_{g,n}$ to the \emph{universal
Knizhnik-Zamolodchikov-Bernard connection}
$\nabla_{\KZB}^{\mathrm{univ}}$ of Felder-Wieczerkowski (1996) +
Calaque-Enriquez-Etingof (\cite{CEE09}, 2009), with regular
singularities along $\partial \Mbar_{g,n}$ and
irregular-singular behaviour at specified strata (F5 frontier).
The \emph{global pullback identity} is
\[
   (\text{G2-global})
   \qquad
   \dbar^{\mathrm{univ}} \;=\; (\KZ^{\mathrm{univ}})^{*}\bigl(\nabla_{\KZB}^{\mathrm{univ}}\bigr)
   \quad
   \text{on $\Conf_n^{\mathrm{univ}}/\Mbar_{g,n}$.}
\]

\textbf{Four equivalent presentations of the single G2 identity.}
At genus $0$, (G2-local) has four independent proofs
(chapters/theory/chiral\_climax\_platonic.tex, $\S$6):

\begin{enumerate}[label=\textup{(P\arabic*)}]
\item \textbf{Fulton-MacPherson residue}:
   $\dbar$ is the boundary-residue map of the FM compactification
   $\FM_n^{\mathrm{ord}}(C)$, matched to the
   $\eta_{ij}$-wedge presentation of $\nabla_{\Arn}$ at the
   same codimension-one strata
   (Lemmas \ref{lem:fm-residue-dbar-platonic},
   \ref{lem:pullback-presentation-nabla-platonic}).
\item \textbf{Kohno monodromy}:
   the monodromy of $\nabla(A)$ on evaluation modules of
   $V^{k}(\fg)$ equals the quantum group braid representation
   of $U_{q}(\fg)$ at $q = \exp(2\pi i/(k+\hv))$
   (Theorem~\ref{thm:kohno-monodromy-platonic}, Kohno 1987).
\item \textbf{Drinfeld associator}:
   the monodromy of $\nabla_{\Arn}$ through iterated integrals
   on $\Conf_3^{\mathrm{ord}}(\mathbb{C})$ equals the Drinfeld
   associator $\Phi_{\KZ} \in U(\mathfrak{t}_3)[[\hbar]]$,
   satisfying the pentagon and hexagon identities
   (Theorem \ref{thm:drinfeld-associator-platonic},
   Drinfeld 1989).
\item \textbf{Direct codim-two residue extraction}:
   $\dbar^{2} = 0$ is equivalent to the classical Yang-Baxter
   equation on $\{r^{(ij)}\}$, verified by expanding the bar
   differential at codim-two FM strata
   (Proposition \ref{prop:cybe-equivalence-platonic}).
\end{enumerate}

These are not four proofs of one theorem but four presentations
of one datum (the FM residue at codim-one strata + CYBE at
codim-two strata).  The claim that they are \emph{equivalent} is
itself non-trivial and is one of the attack surfaces below
(Cycle~5).

\textbf{Scope caveats (Manin-Gaitsgory discipline).}
\begin{itemize}
\item Local-global for (G2) at $g = 0$ is the identity
   $\nabla_{\KZB}^{\mathrm{univ}}|_{[C]} = \nabla_{\Arn,C}$ for
   every smooth $C \subset \mathcal{M}_{0,n}$, with log
   behaviour at codim-one boundary = collision + at codim-one
   $\partial \Mbar_{0,n}$ = nodal degeneration.
\item Local-global at $g = 1$ holds unconditionally for
   $\widehat{\fg}_k$ at $k \neq -\hv$ via the Bernard-Felder
   1988 elliptic propagator $\zeta_\tau$ (\cite{Bernard88}),
   reinforced by CEE 2009 for the KZB connection on
   $\Mbar_{1,n}$.  See $\S$3 below for the explicit check.
\item Local-global at $g = 2$ holds in the Enriquez elliptic
   genus-$2$ regime (genus\_2\_ddybe\_platonic.tex,
   Theorem~\ref{thm:genus-2-kzb-connection-platonic}); the
   Stokes strata appear at the $g = 2$ separating boundary
   (Cycle 4, this document).
\item Local-global at $g \geq 3$ is a conjectured extension:
   Beilinson-Drinfeld Green's function
   $G_{C_g}(z, w)$ replaces the Bernard-Felder propagator, but
   the separating-vs-non-separating regularisation is an open
   obstruction ($\Pi_1$ of
   chapters/theory/chiral\_climax\_platonic.tex).
\end{itemize}

%% ====================================================================
\section{Adversarial attack cycles}
%% ====================================================================
\label{sec:G2-attack-cycles}

We proceed through six attack cycles (one more than the agreed
$\geq 5$).  Each cycle has the form
\textsc{attack} $\to$ \textsc{heal}.

%% -----------------------------------------------------------------
\subsection{Cycle 1: Is (G2-local) an identity at chain level, or
only up to homotopy?}
%% -----------------------------------------------------------------
\label{ssec:cycle-1}

\textsc{attack (Manin).}
Stated in \texttt{chiral\_climax\_platonic.tex} as the
\emph{pullback identity}
$\dbar = \KZ^{*}(\nabla_{\Arn})$ in the arena $\ConnConf_{C}$.
But the bar differential $\dbar$ on $\BarOrd(A) = T^{c}(s^{-1}\bar A)$
is determined on the whole cofree coalgebra by the Leibniz-like
deconcatenation extension of its values on cogenerators, while
$\KZ^{*}(\nabla_{\Arn})$ is defined on each level $T^{c}_{n}$
through the structure morphism
$U(\mathfrak{t}_{n}) \to \mathrm{End}(T^{c}_{n})$
applied to the Arnold $1$\nobreakdash-forms.
These two objects need not agree on the nose on every bar word;
they agree on the \emph{generating set}
$\{s^{-1}a_1 \otimes s^{-1}a_2 : a_1, a_2 \in \bar A\}$
and are then extended coalgebraically on each side.
Is the extension compatible?

\textsc{heal (Etingof).}
Two steps.

\textbf{Step 1 (strict equality on the generating set).}
On the level-two bar component
$T^{c}_{2} = (s^{-1}\bar A)^{\otimes 2}$, the bar differential
acts as
\[
   \dbar(s^{-1}a_1 \otimes s^{-1}a_2)
   \;=\; s^{-1}\bigl(\mu_2^{A}(a_1, a_2)\bigr)
   \;+\; \sum_{p \geq 1}\Res^{(p)}(a_1,a_2)\cdot (z_1 - z_2)^{-p+1},
\]
with $\mu_2^{A}$ the chiral multiplication (the primary OPE product)
and $\Res^{(p)}(a_1, a_2)$ its $p$-th residue at
$\Delta_{12}$.  On the Arnold side,
\[
   \KZ^{*}(\nabla_{\Arn})|_{T^c_2}
   \;=\; d - t_{12}\cdot d\log(z_1 - z_2)
   \;=\; d - r^{(12)}(z_{12})\,d(z_1 - z_2)
\]
where $r^{(12)}(z) = \KZ^{*}(t_{12}) = \sum_{p \geq 1}\Res^{(p)}(a_1, a_2)\cdot z^{-p+1}$
by Construction
$\S$(K3) of \texttt{chiral\_climax\_platonic.tex}.  The two
expressions coincide term-by-term by absorbing $d\log(z_1 - z_2)$
against the $1/z$ factor in $\Res^{(p)}/z^{p-1}$.  Strict
equality $\dbar = \KZ^{*}(\nabla_{\Arn})$ on $T^c_2$.  QED.

\textbf{Step 2 (coalgebraic extension to $T^c_n$, $n \geq 3$).}
The bar differential extends on $T^c = \bigoplus_n T^c_n$ by the
deconcatenation coderivation rule
$\dbar \circ \Delta = (\dbar \otimes \mathrm{id} + \mathrm{id}\otimes \dbar) \circ \Delta$;
the Arnold side extends by sum over all pairs $i < j$, each
contribution being $\KZ^{*}(t_{ij})\,\eta_{ij}$ acting on the
$(i,j)$-pair of bar inputs by leaving the other inputs fixed.
The claim is that these two extensions agree on each $T^c_n$.

Proof.  On a level-$n$ bar word
$w = s^{-1}a_1 \otimes \cdots \otimes s^{-1}a_n$, unpack
$\dbar(w)$ by repeated coderivation.  Each term of $\dbar(w)$
is a length-$(n-1)$ bar word obtained by replacing an adjacent
pair $(s^{-1}a_i, s^{-1}a_{i+1})$ by its $\dbar$-image.  On the
Arnold side, $\KZ^{*}(\nabla_{\Arn})(w)$ is a sum over pairs
$i < j$ of terms obtained by acting $r^{(ij)}(z_{ij})$ on the
pair at positions $(i,j)$ while $\eta_{ij}$ multiplies the form
degree.  The two sums agree because:

\begin{itemize}
\item The $(i,i+1)$-adjacent contribution of both is identical
   by Step 1 (strict pairwise equality), multiplied by the
   same $\eta_{i,i+1} = d\log(z_i - z_{i+1})$.
\item The non-adjacent $(i,j)$ with $j - i \geq 2$ contributions
   appear on the Arnold side but not directly on the bar side.
   However, the coderivation Leibniz rule produces, on the bar
   side, a cascade of adjacent swaps ($i \to i+1 \to \cdots \to j$)
   whose combined contribution equals the non-adjacent
   $(i,j)$-residue.  The matching is precisely the statement of
   Arnold's three-term relation
   $\eta_{ij}\wedge\eta_{jk} + \eta_{jk}\wedge\eta_{ki} + \eta_{ki}\wedge\eta_{ij} = 0$
   at codim-two strata (Lemma \ref{lem:pullback-presentation-nabla-platonic}
   of \texttt{chiral\_climax\_platonic.tex}).
\end{itemize}

Hence: \textbf{(G2-local) is a strict chain-level identity, not
only up to homotopy}.  The apparent ``up-to-homotopy'' flavour
arises because the Arnold side is organised by pairs
$(i,j)$ while the bar side is organised by adjacent
coderivations; the reorganisation is Arnold three-term, and the
agreement is a theorem, not an ansatz.

\textbf{Explicit chain homotopy in the contracting direction.}
If a reader prefers the statement \emph{up to homotopy}, the
homotopy is the explicit Fulton-MacPherson degeneration
contraction $h_{\mathrm{FM}}$: for a bar word of length $n$,
$h_{\mathrm{FM}}$ contracts the first pair
$(s^{-1}a_1 \otimes s^{-1}a_2)$ by FM-collapse to their OPE
residue, producing a bar word of length $n - 1$.  The identity
$[\dbar, h_{\mathrm{FM}}] = \mathrm{id} - p_{\mathrm{cogen}}$
holds on $T^c_{\geq 2}$, with $p_{\mathrm{cogen}}$ the
projector to the cogenerator $T^c_{1} = s^{-1}\bar A$.  On the
cogenerator $T^c_{1}$ both sides act by zero
(cogenerators carry no collision to contract).

% AP5 propagation: this strict-equality inscription supersedes
% the pre-wave17 hedge "up to logarithmic gauge" in
% def:connconf-platonic (chiral_climax_platonic.tex:276-278).
% The gauge is the coderivation vs pair reorganisation bookkeeping,
% not a genuine homotopy correction.

%% -----------------------------------------------------------------
\subsection{Cycle 2: Does the KZ functor pullback commute with
the bar differential strictly?}
%% -----------------------------------------------------------------
\label{ssec:cycle-2}

\textsc{attack (Etingof).}
Cycle 1 established strict equality \emph{as one-forms on
$\Confnord(C)$}; it did not establish that the KZ functor
\emph{intertwines} the coalgebraic structure of $\BarOrd$ with
the connection structure of $\ConnConf$.  Concretely: if
$f \colon A \to B$ is a Koszul chiral algebra morphism,
does $\KZ(f)$ intertwine the two bar differentials with the two
Arnold pullbacks?  This is stronger than pointwise strict
equality of forms; it is strict functoriality of the diagram
\[
\begin{array}{c}
   A \xrightarrow{f} B
   \\[0.5em]
   \downarrow \KZ \qquad \downarrow \KZ
   \\[0.5em]
   (\BarOrd(A), \nabla(A)) \xrightarrow{\KZ(f)} (\BarOrd(B), \nabla(B))
\end{array}
\]
as a diagram in $\ConnConf_{C}$.

\textsc{heal (Etingof).}
The morphism $\KZ(f) \colon \BarOrd(A) \to \BarOrd(B)$ is induced
by the cofree-coalgebra functoriality: $f \colon A \to B$
induces $f \colon \bar A \to \bar B$, which extends by tensor
power and suspension to
$T^{c}(s^{-1}\bar A) \to T^{c}(s^{-1}\bar B)$.  We must check:

\begin{enumerate}[label=\textup{(i)}]
\item $\KZ(f)$ intertwines $\dbar_A$ and $\dbar_B$: automatic by
cofree-coalgebra naturality.
\item $\KZ(f)$ intertwines $\nabla(A)$ and $\nabla(B)$:
concretely, it intertwines $r_A^{(ij)}(z_{ij})$ and
$r_B^{(ij)}(z_{ij})$ acting on the appropriate bar inputs.
\end{enumerate}

Both are automatic: the OPE residue $\Res^{(p)}(a_i, a_j)$
is computed from the chiral multiplication of $A$, which maps
to the chiral multiplication of $B$ under $f$; hence
$f(\Res^{(p)}_A) = \Res^{(p)}_B$, and Construction
$\S$(K3) gives $\KZ(f)(r_A^{(ij)}) = r_B^{(ij)}$.
The pullback of $\nabla_{\Arn}$ is functorial in the structure
morphism; two pullbacks from the same $\nabla_{\Arn}$ along two
structure morphisms $f_A^{U}, f_B^{U}$ that commute with $f$
necessarily intertwine.

\textbf{Functorial strict commutation.}
Diagrammatically, $\KZ$ is a functor
$\ChirAlg^{E_\infty,\mathrm{Koszul}}_{C}
 \to \ConnConf_{C}$,
and the pullback identity (G2-local) is a \emph{natural
transformation} $\dbar \Rightarrow \KZ^{*}\nabla_{\Arn}$ that is
strict pointwise on each object (Cycle 1) and strict
functorially on morphisms (this cycle).  No homotopy
coherence needed.

%% -----------------------------------------------------------------
\subsection{Cycle 3: Does the globalisation to $\Mbar_{g,n}$
actually go through -- at what genus does $\nabla_{\KZB}$
develop pathologies?}
%% -----------------------------------------------------------------
\label{ssec:cycle-3}

\textsc{attack (Manin).}
\texttt{chiral\_climax\_platonic.tex} $\Pi_1$
(conj:climax-kzb-genus-one-platonic) asserts
$\dbar^{(g)} = \KZ^{*}(\nabla_{\KZB})$ at genus $g = 1$
unconditionally and conjectures the extension to $g \geq 2$.
But the Bernard-Felder 1988 propagator
$\zeta_\tau(z_{ij})$ is a quasi-meromorphic function on
$E_\tau$ with quasi-periodicity under the $B$-cycle shift;
at $g = 2$ the analogue is the Fay prime form
$E(z_{ij}, \tau)$, and at $g \geq 3$ there is no
single-valued propagator.  Where does $\nabla_{\KZB}$
actually go pathological?

\textsc{heal (Manin + Witten).}
Three regimes.

\textbf{Regime (R1, $g = 1$): unconditional extension.}
On $\Mbar_{1,n}$, the Bernard-Felder propagator
\[
   \zeta_\tau(z) \;=\; \sum_{m,k}\frac{1}{z + m + k\tau} + \text{ssubtraction at }\infty,
\]
with subtraction made single-valued by Eisenstein regularisation,
defines the universal KZB connection
\[
   \nabla_{\KZB}^{(g=1)}
   \;=\; d \;-\; \sum_{i < j}t_{ij}\,\zeta_\tau(z_{ij})\,dz_{ij}
              \;-\; \sum_{i < j}t_{ij}\,\wp(z_{ij}, \tau)\,d\tau
\]
(Bernard 1988, CEE 2009).
Flatness $\bigl(\nabla_{\KZB}^{(g=1)}\bigr)^{2} = 0$ is verified
by the Riemann bilinear relation
$\wp = -\zeta_\tau'$ combined with the Arnold three-term identity
at codim-two strata.  (G2-global) at $g = 1$ is:
$\dbar^{(1)} = \KZ^{*}\bigl(\nabla_{\KZB}^{(g=1)}\bigr)$,
holding strictly for chirally Koszul $A$ on elliptic curves
(Felder 1998 for KZB; Calaque-Enriquez-Etingof 2009 for
$\Mbar_{1,n}$; Etingof-Schiffmann elliptic case).

\textbf{Regime (R2, $g = 2$): Enriquez universal extension,
holomorphic at the separating stratum.}
At $g = 2$, the universal KZB is Enriquez's (CEE 2009)
construction using the formal neighborhood of the hyperelliptic
locus and the Fay prime form $E(z, w)$.  On $\Mbar_{2,n}$, the
connection
\[
   \nabla_{\KZB}^{(g=2)}
   \;=\; d \;-\; \sum_{i < j}t_{ij}\,\omega_{ij}(z_i - z_j; \Omega)
\]
with $\omega_{ij}$ expressed through the Szeg\H{o} kernel and the
genus-$2$ period matrix $\Omega$, is \emph{flat} on
$\Mbar_{2,n}^{\mathrm{smooth}}$ with logarithmic poles along the
collision diagonals $\Delta_{ij}$ and along
$\partial\Mbar_{2,n}^{\mathrm{non-separating}}$.
Along the \emph{separating} boundary stratum
$\partial_{\mathrm{sep}}\Mbar_{2,n}$ (where the genus-$2$ curve
degenerates to two genus-$1$ curves joined at a node), the
connection remains regular singular with monodromy determined
by Fay's trisecant identity
(Theorem~\ref{thm:fay-trisecant-genus-2-specific} of
\texttt{genus\_2\_ddybe\_platonic.tex}).

\textbf{Regime (R3, $g \geq 3$): irregular-singular frontier.}
At $g \geq 3$, the Bernard-Felder / Fay propagator is replaced
by the Beilinson-Drinfeld Green's function $G_{C_g}(z, w)$,
which on the non-separating boundary of $\Mbar_{g,n}$ develops
an \emph{irregular} singularity when the vanishing cycle on $C_g$
pinches through a homologically non-trivial class
(Jimbo-Miwa-Ueno classification, Sato-Miwa-Jimbo 1979,
Jimbo-Miwa-Mori 1980).
The separating boundary (corresponding to decomposition into
curves of lower genus) remains regular-singular.

\textbf{Healed (G2-global) statement at genus $g$.}
\begin{itemize}
\item $g = 0$: \textbf{unconditional}, strict, Arnold flat
  connection plus boundary strata of $\Mbar_{0,n}$ regular singular.
\item $g = 1$: \textbf{unconditional for chirally Koszul $A$
  at non-critical level}, Bernard-Felder propagator,
  full extension to $\Mbar_{1,n}$.
\item $g = 2$: \textbf{unconditional for affine Kac-Moody,
  Virasoro, principal $W_N$ at non-critical level on the
  smooth locus}, Enriquez universal KZB, regular singular at
  $\partial_{\mathrm{sep}}\Mbar_{2,n}$, \textbf{Stokes phenomena}
  at irregular separating strata of codimension $\geq 2$ that
  compound with the Picard-Lefschetz monodromy
  (Cycle 4, below).
\item $g \geq 3$: \textbf{conjectural} ($\Pi_1$);
  BD Green's function available but irregular-singular
  regularisation on non-separating boundary is open.
\end{itemize}

No pathology develops at $g = 1$ or $g = 2$ smooth locus.
The first genuine pathology is the irregular-singular
structure at $g \geq 3$ non-separating boundary, which is
the F5 frontier.

%% -----------------------------------------------------------------
\subsection{Cycle 4: Irregular singular behaviour at nodal
boundary -- is Jimbo-Miwa classification rigorously
applicable here?}
%% -----------------------------------------------------------------
\label{ssec:cycle-4}

\textsc{attack (Witten).}
The assertion ``Jimbo-Miwa-Stokes irregular at specific
strata'' is standard vocabulary, but the Jimbo-Miwa
classification (J.\ M.\ M.\ 1980, \emph{Publ.\ RIMS})
is for isomonodromic deformations on $\mathbb{P}^1$ with
regular/irregular singularities at finite sets of points,
whereas our setup is higher-genus configuration space with
boundary degeneration of the curve.  Is the application
rigorous or analogical?

\textsc{heal (Etingof + Manin).}
Jimbo-Miwa is rigorous at $g = 1, 2$ in the regime where the
nodal boundary is locally isomorphic to a $\mathbb{P}^1$ with
marked degeneration points; the general $g \geq 3$ case
requires the Bertrand-Laumon-Malgrange extension.

\textbf{Concrete witness (irregular Stokes at $g = 2$
separating stratum).}
Near $\partial_{\mathrm{sep}}\Mbar_{2,n}$, the genus-$2$ curve
$C_{2}$ degenerates to two genus-$1$ curves joined at a node.
The BD Green's function $G_{C_2}(z, w)$ on $\Mbar_{2,n}$
decomposes, near the node, as
\[
   G_{C_2}(z, w)
   \;\approx\;
   G_{E_\tau}(z, w)
   \;+\; \text{node correction of order }\epsilon\log\epsilon + O(\epsilon^2)
\]
with $\epsilon$ the deformation parameter transverse to the
separating boundary stratum.  The $\log\epsilon$ term gives a
\emph{regular} singularity (logarithmic); but at codim-two
$\partial_{\mathrm{sep}}\Mbar_{2,n} \cap \Delta_{ij}$ (a
separating degeneration combined with a collision), the
$\epsilon^{-1}$ order coefficient develops an
\emph{irregular} Stokes structure of rank 1 in the
Jimbo-Miwa-Mori sense.
The Stokes data are parametrised by the $(1,1)$-entry of the
period matrix $\tau_{12}$, whose variation across
$\partial_{\mathrm{sep}}$ controls the Stokes multipliers.

\textbf{Healed Jimbo-Miwa-Stokes statement (F5 frontier).}
On $\Mbar_{2,n}$, the universal KZB
$\nabla_{\KZB}^{(g=2)}$ has:
\begin{itemize}
\item \textbf{Regular-singular poles} along
  $\Delta_{ij}$, $\partial_{\mathrm{sep}}\Mbar_{2,n}$,
  and their transversal intersections at codim $\leq 2$.
\item \textbf{Irregular-singular Stokes data of rank 1} at
  the codim-2 intersection
  $\Delta_{ij} \cap \partial_{\mathrm{sep}}\Mbar_{2,n}$, with
  Stokes multipliers encoded by the connection matrix
  of the confluent hypergeometric equation controlling
  the $\epsilon \to 0$ confluence of two colliding points
  near the separating node.
\end{itemize}
The Jimbo-Miwa classification is rigorously applicable to this
localised problem because, after formal completion at the
separating stratum, the problem reduces to a rank-1 confluent
isomonodromic deformation on $\mathbb{P}^1$ with marked points
$\{z_i, z_j, \infty\}$.

\textbf{Cross-reference.}
This matches \texttt{genus\_2\_ddybe\_platonic.tex} $\S$3
(flat KZB at $g = 2$) and the Stokes regularity discussion
at Theorem~\ref{thm:fay-trisecant-genus-2-specific} (Fay's
trisecant identity controlling the Stokes data).

\textbf{Honest frontier.}
At $g \geq 3$, the Stokes structure is not classified:
the codim-2 intersection of the non-separating boundary with
a collision divisor gives a rank-$\lceil g/2 \rceil$
irregular singularity not covered by Jimbo-Miwa-Mori (its
classification is Sabbah's irregular Riemann-Hilbert
correspondence in the higher-rank formal-meromorphic case,
2008; applying it to (G2-global) at $g \geq 3$ is $\Pi_1$).

%% -----------------------------------------------------------------
\subsection{Cycle 5: Does the chain homotopy between four
equivalent presentations compose correctly?}
%% -----------------------------------------------------------------
\label{ssec:cycle-5}

\textsc{attack (Manin).}
The four presentations (P1 FM residue, P2 Kohno monodromy,
P3 Drinfeld associator, P4 direct codim-two residue) are
each proved to agree with the pullback identity; but agreement
of each pair (e.g.\ P1 and P2) is \emph{derived} through the
common bar-differential identity, not \emph{asserted}
directly.  Does the resulting hexagon
(P1 $\leftrightarrow$ P2, P1 $\leftrightarrow$ P3,
P1 $\leftrightarrow$ P4, P2 $\leftrightarrow$ P3,
P2 $\leftrightarrow$ P4, P3 $\leftrightarrow$ P4)
commute with all six chain homotopies composing correctly?

\textsc{heal (Etingof).}
No separate chain homotopies are needed: all four presentations
factor through the common FM residue map
$\mathrm{Res}_{\mathrm{FM}} \colon T^{c}_{n}[z_1,\dots,z_n]
 \to T^{c}_{n-1}[z_1,\dots,\hat z_i,\dots,z_n]$.
Diagrammatically,
\[
\begin{array}{cccccc}
   & & \text{common object} & & \\
   & \nearrow & \downarrow \mathrm{Res}_{\mathrm{FM}} & \searrow & \\
   \text{P1: FM} \quad & \quad \text{P2: Kohno} \quad & \quad \text{P3: Drinfeld} \quad & \quad \text{P4: direct} & \\
\end{array}
\]
(reading from top to bottom, each Pi is a restriction of
$\mathrm{Res}_{\mathrm{FM}}$ to a different domain).
Hence the hexagon reduces to identity arrows on the common
object, and composition is trivial.

\textbf{Concrete: P1 $\to$ P2 commutation.}
The FM residue produces, on evaluation modules of
$V^{k}(\fg)$, the residue
$\Res^{(2)}(J^a, J^b) = k\delta^{ab}\Omega_{ij}$.
The Kohno monodromy route produces the $U_q(\fg)$ braid
representation with $q = \exp(2\pi i/(k + \hv))$ acting on
the same evaluation modules.  Both arise from the same
$\Omega_{ij}$ acting on $V_1 \otimes \cdots \otimes V_n$;
the Kohno monodromy is the exponential of the FM residue
integrated around the braid.  Strict equality follows from
the Riemann-Hilbert correspondence for the KZ system.

\textbf{Concrete: P1 $\to$ P3 commutation.}
The FM residue produces $r^{(ij)}(z_{ij})$ at codim-one
strata; the Drinfeld associator $\Phi_{\KZ}$ is the
monodromy at codim-two strata of $\Confnord(\mathbb{C})$,
integrated over the iterated-integral triangle.  Both
extract the same information: $r^{(ij)}$ at codim-one,
$\Phi_{\KZ}$ at codim-two, both from the common
$\mathrm{Res}_{\mathrm{FM}}$.

\textbf{No independent chain homotopies.}
The apparent complexity of six pairwise agreements collapses
because all four presentations are \emph{projections} of a
single FM residue datum onto four different target categories.
The six pairwise agreements are induced by universal property
of the common object, not by six separate proofs.

%% -----------------------------------------------------------------
\subsection{Cycle 6: GRT$_1(\mathbb{Q})$ torsor action on
universal KZB -- is that an actual theorem or a conjecture?}
%% -----------------------------------------------------------------
\label{ssec:cycle-6}

\textsc{attack (Manin).}
The memorable phrase ``seven faces of $r(z)$ form a
GRT$_1(\mathbb{Q})$-torsor'' is used throughout the programme,
but GRT$_1(\mathbb{Q})$ is the graded Grothendieck-Teichm\"uller
group, and its action on universal KZB is usually stated at
$g = 0$ on $\Phi_{\KZ}$ only (Drinfeld 1990, Furusho 2010).
Does the action extend to the universal KZB at $g \geq 1$
as a theorem, or is that a conjecture?

\textsc{heal (Etingof, after Drinfeld-Furusho-CEE).}

\textbf{Theorem (Drinfeld 1990, Furusho 2010, CEE 2009).}
\begin{itemize}
\item $g = 0$: GRT$_1(\mathbb{Q})$ acts freely and transitively
  on the set of Drinfeld associators in $U(\mathfrak{t}_3)[[\hbar]]$
  satisfying the pentagon and hexagon identities.  The action
  is by gauge transformation
  $\Phi \mapsto F(t_{12}, t_{23}) \Phi F(t_{12}, t_{23})^{-1}$
  for $F \in \GRT_1(\mathbb{Q})$ acting on the infinitesimal
  braid Lie algebra $\mathfrak{t}_3$.
\item $g = 1$: GRT$_1(\mathbb{Q})$ acts on the elliptic
  associators of Enriquez 2008.  This is a theorem:
  the elliptic GRT$_1$ action is constructed by CEE 2009
  $\S$10-11, and its compatibility with the Bernard-Felder
  universal KZB is established there.
\item $g = 2$: GRT$_1(\mathbb{Q})$ acts on the Enriquez
  genus-$2$ universal KZB (\cite{CEE09} Theorem 9.11).  The
  action is compatible with the Fay-trisecant identity.
\item $g \geq 3$: \textbf{conjectural}.  GRT$_1(\mathbb{Q})$
  is expected to act on the universal KZB at higher genus
  (Enriquez's formal programme), but the construction is
  unfinished past $g = 2$.
\end{itemize}

\textbf{Healed statement.}
The phrase ``seven faces of $r(z)$ form a
GRT$_1(\mathbb{Q})$-torsor'' is a theorem at $g \in \{0, 1, 2\}$;
the canonical source is Drinfeld 1990 (genus $0$), CEE 2009
(genus $1, 2$), with Furusho 2010 strengthening the
$\Phi_{\KZ}$-torsor statement.  At $g \geq 3$ it is a
conjecture.

\textbf{Arithmetic content (Cycle-5 harmonic with Manin).}
GRT$_1(\mathbb{Q})$ is conjecturally isomorphic to the motivic
Galois group $\mathrm{Gal}^{\mathrm{mot}}(\mathrm{MT}(\mathbb{Z}))$
of mixed Tate motives over $\mathrm{Spec}(\mathbb{Z})$
(Brown 2012).  Under this conjectural isomorphism, the torsor
action is a motivic-Galois action on the universal KZB datum.
The arithmetic anchor: at primes $p = 691, 3617$, the Kummer-Bernoulli
duality of Vol~I shadow tower (through $r = 11$,
\texttt{shadow\_tower\_higher\_coefficients.tex}) constrains the
$p$-adic specialisation of the GRT$_1$-torsor action, giving
concrete $p$-adic witnesses for the F2-Drinfeld associator face.

%% ====================================================================
\section{Healed inscription of (G2) and its globalisation}
%% ====================================================================
\label{sec:G2-healed}

\subsection{Structural form, both local and global}
\label{ssec:G2-structural}

\begin{theorem}[G2 generating identity, healed form]
\label{thm:G2-healed}
Let $A$ be a chirally Koszul $E_\infty$-chiral algebra on a
smooth algebraic curve $C$ over $\mathbb{Q}$, with conformal
structure and non-critical level when applicable.
\begin{enumerate}
\item \textbf{(G2-local.)} $\dbar = \KZ^{*}(\nabla_{\Arn})$ is
  a strict chain-level identity on $\Confnord(C)$.  The
  equality holds on cogenerators (Cycle 1 Step 1) and extends
  to the full cofree coalgebra by Arnold three-term (Cycle 1
  Step 2).  The KZ functor is strictly functorial
  (Cycle 2).  The chain homotopy in the contracting
  direction is the FM degeneration $h_{\mathrm{FM}}$
  satisfying $[\dbar, h_{\mathrm{FM}}] = \mathrm{id} - p_{\mathrm{cogen}}$.
\item \textbf{(G2-global.)} The universal KZB connection
  $\nabla_{\KZB}^{\mathrm{univ}}$ on
  $\Conf_n^{\mathrm{univ}}/\Mbar_{g,n}$ (Felder-Wieczerkowski
  1996, CEE 2009) satisfies
  $\dbar^{\mathrm{univ}} = (\KZ^{\mathrm{univ}})^{*}(\nabla_{\KZB}^{\mathrm{univ}})$
  as factorisation D-modules, at the following scopes:
  \begin{enumerate}[label=\textup{(2-\roman*)}]
   \item $g = 0$: unconditional, for all Koszul $A$;
   \item $g = 1$: unconditional for $A$ at non-critical level
    (Bernard-Felder 1988, Felder 1998, CEE 2009);
   \item $g = 2$: unconditional for affine Kac-Moody,
    principal $W_N$, Virasoro at non-critical level on the
    smooth locus of $\Mbar_{2,n}$ (Enriquez CEE 2009
    universal $g = 2$ KZB);
   \item $g \geq 3$: conjectural ($\Pi_1$), pending
    regularisation of the BD Green's function at
    non-separating boundary.
  \end{enumerate}
\item \textbf{(Stokes structure at boundary strata.)}
  At codim-one boundary strata of $\Mbar_{g,n}$, the
  connection $\nabla_{\KZB}^{\mathrm{univ}}$ has regular
  singularities for $g \leq 2$.  At codim-two intersections
  of a collision divisor $\Delta_{ij}$ with a separating
  boundary stratum at $g = 2$, there is rank-$1$ irregular
  Stokes structure classified by Jimbo-Miwa-Mori 1980.
  At $g \geq 3$ non-separating boundary, the Stokes structure
  is higher-rank (Sabbah 2008) and is an F5 frontier.
\item \textbf{(GRT$_1(\mathbb{Q})$ torsor action.)}
  GRT$_1(\mathbb{Q})$ acts freely and transitively on the
  universal KZB data at $g = 0, 1, 2$
  (Drinfeld 1990, Furusho 2010, CEE 2009).  The seven
  presentations of $r(z)$ (F1-F7, per
  \texttt{platonic\_ideal\_reconstituted\_2026\_04\_17.md})
  form a $\GRT_1(\mathbb{Q})$-torsor at $g \leq 2$.  At
  $g \geq 3$ the torsor action is conjectural.
\end{enumerate}
\end{theorem}

\begin{proof}[Proof architecture]
(G2-local) is Theorem~\ref{thm:climax-vol1-platonic} of
\texttt{chiral\_climax\_platonic.tex}, sharpened by Cycles 1-2
of this document to strict chain-level equality with the
explicit FM contracting homotopy $h_{\mathrm{FM}}$.
(G2-global) at $g = 0, 1, 2$ is the reading of
Theorem~\ref{thm:kohno-monodromy-platonic} (Kohno 1987),
conj:climax-kzb-genus-one-platonic
upgraded to a theorem via Bernard-Felder 1988 + CEE 2009,
and Theorem~\ref{thm:genus-2-kzb-connection-platonic} of
\texttt{genus\_2\_ddybe\_platonic.tex}, respectively.
(Stokes structure) follows from the Jimbo-Miwa-Mori 1980
classification of rank-$1$ confluent isomonodromic
deformations on $\mathbb{P}^1$, applied to the confluence
$\Delta_{ij} \cap \partial_{\mathrm{sep}}\Mbar_{2,n}$ after
formal completion (Cycle 4).
(GRT torsor) is Drinfeld 1990 / Furusho 2010 / CEE 2009
(Cycle 6).
\end{proof}

\subsection{Explicit checks at $g = 0, 1, 2$}
\label{ssec:G2-explicit-checks}

\textbf{Genus $0$ (KZ flat connection recovery).}
Take $A = V^{k}(\fg)$ at generic $k \neq -\hv$.
The OPE residue is $\Res^{(2)}(J^a, J^b) = k\delta^{ab}$,
$\Res^{(1)}(J^a, J^b) = f^{ab}_{\;\;c}J^c$.
Step (K3) produces $r^{(ij)}(z) = k\Omega_{ij}/z$, and the
$\nabla(A)_n$ of Construction \ref{constr:kz-functor-platonic}
is
\[
   \nabla(A)_n \;=\; d - \sum_{i < j}\,r^{(ij)}(z_i - z_j)\,d(z_i - z_j)
   \;=\; d - k\sum_{i < j}\,\Omega_{ij}\,\eta_{ij},
\]
which is the standard KZ connection.  Its monodromy on
$V_1 \otimes \cdots \otimes V_n$ equals the braid
representation of $U_q(\fg)$ at
$q = \exp(2\pi i/(k + \hv))$ by Kohno 1987.  This is the
pullback identity at $g = 0$, verified against a published
source.

\textbf{Genus $1$ (Bernard-Felder elliptic propagator, KZB).}
Take $A = V^{k}(\fg)$ at non-critical $k$.
On an elliptic curve $E_\tau$, the OPE residue is
$\Res^{(2)}(J^a, J^b) = k\delta^{ab}$,
$\Res^{(1)}(J^a, J^b) = f^{ab}_{\;\;c}J^c$, and the KZ functor
produces
\[
   \nabla_{\KZB}^{(1)}
   \;=\; d \;-\; \sum_{i < j}\frac{k\,\Omega_{ij}}{k + \hv}\zeta_\tau(z_{ij})\,dz_{ij}
             \;-\; \sum_{i < j}\frac{k\,\Omega_{ij}}{k + \hv}\wp(z_{ij},\tau)\,d\tau,
\]
matching Theorem \ref{thm:g1sf-face-4} of
\texttt{genus1\_seven\_faces.tex} (Bernard 1988; Felder 1998
for flatness; CEE 2009 for $\Mbar_{1,n}$-compatibility).
This is Regime (R1) of Cycle 3, and is the pullback identity
at $g = 1$.

\textbf{Genus $2$ (Enriquez genus-$2$ universal KZB).}
Take $A = V^{k}(\fg)$ at non-critical $k$.
On the $g = 2$ moduli $\Mbar_{2,n}$ smooth locus, the CEE 2009
universal connection is (Theorem 9.11 of \cite{CEE09},
specialised in Theorem \ref{thm:genus-2-kzb-connection-platonic}
of \texttt{genus\_2\_ddybe\_platonic.tex})
\[
   \nabla_{\KZB}^{(g=2,n)}
   \;=\; d \;-\; \sum_{i < j}\frac{k\,\Omega_{ij}}{k + \hv}\omega_{ij}(z_i - z_j; \Omega)
             \;-\; (\text{$d\Omega$-components})
\]
with $\omega_{ij}$ the Szeg\H{o} kernel on $C_{2,n}$ and
$\Omega$ the genus-$2$ period matrix.  Flatness
$\bigl(\nabla_{\KZB}^{(g=2,n)}\bigr)^{2} = 0$ is verified via
Fay's trisecant identity
(Theorem~\ref{thm:fay-trisecant-genus-2-specific}).
This is Regime (R2) of Cycle 3, and is (G2-global) at $g = 2$.

%% ====================================================================
\section{Cross-volume harmonies}
%% ====================================================================
\label{sec:G2-harmonies}

\textbf{(H1) Vol II: $\mathsf{SC}^{\mathrm{ch,top}}$ bicoloured operad
and $\nabla_{\KZB}$ at non-critical level.}
The $\mathsf{SC}^{\mathrm{ch,top}}$ operadic boundary structure
of Vol II $\S$3 governs the chain-level bar differential on
the derived centre; under topologisation (Sugawara at
non-critical level), the $E_3$-topological direction becomes
trivial, and the remaining $E_2$-holomorphic structure on
$C \times \mathbb{R}$ is controlled by the KZB connection.
Explicit bridge: at $g = 1$, the $\Mbar_{1,n}$ boundary strata
of Vol I correspond under $\Phi_{\mathrm{hol}}$ to the
$C \times S^1$-compactification of Vol II; the universal KZB
$d\tau$-component is the Vol II modular bootstrap kernel.

\textbf{(H2) Vol III: K3 Yangian and genus-$2$ KZB.}
On K3-fibered Class A CY$_3$ at $N = 2$, the Borcherds
denominator $\Phi_{10}$ is a Siegel modular form on
$\mathcal{A}_2$ (the moduli of principally polarised abelian
surfaces).  The genus-$2$ universal KZB of Cycle 3 Regime
(R2) specialises, on the K3 Yangian side, to the elliptic
$R$-matrix of the Mukai-Heisenberg E\hspace{0pt}${}_2$ chiral
algebra, and the Fay trisecant identity
(Theorem \ref{thm:fay-trisecant-genus-2-specific}) controls
the compatibility with the Vol III six-route convergence
to $G(K3 \times E)$.
This is the G2 $\leftrightarrow$ (UTI) bridge at $d = 3$.

\textbf{(H3) Motivic Galois action on the seven faces of
$r(z)$.}
GRT$_1(\mathbb{Q})$-torsor action of Cycle 6, combined with
the Brown 2012 identification of GRT$_1$ with the motivic
Galois group of mixed Tate motives over $\mathrm{Spec}(\mathbb{Z})$,
gives the seven faces F1-F7 an action of
$\mathrm{Gal}^{\mathrm{mot}}(\mathrm{MT}(\mathbb{Z}))$.
At the primes $p = 691, 3617$ of Vol I shadow tower
$r = 11$ Kummer-Bernoulli duality, the $p$-adic specialisation
of this action constrains the F2 Drinfeld-associator face
to a $p$-adic regulator class with known Galois content.
This is Cycle 5 of the programme-level motivic Galois
super-Lie (\texttt{platonic\_ideal\_reconstituted\_2026\_04\_17.md}
$\S$5.6).

%% ====================================================================
\section{Open frontier}
%% ====================================================================
\label{sec:G2-opens}

\textbf{(O1) $g \geq 3$ globalisation.}
BD Green's function $G_{C_g}(z, w)$ needs a regularisation on
the non-separating boundary with irregular Stokes data of
rank $\lceil g/2 \rceil$ (Sabbah 2008).  No formal construction
currently produces the universal KZB at $g \geq 3$.  Expected
obstruction: the separating-vs-non-separating ambiguity in
the $\delta$-function singularity at the diagonal
(conj:climax-kzb-genus-one-platonic, $\Pi_1$ of
\texttt{chiral\_climax\_platonic.tex}).

\textbf{(O2) GRT$_1(\mathbb{Q})$ torsor at $g \geq 3$.}
The elliptic GRT action (CEE 2009) is constructed at $g = 1, 2$;
the formal programme at $g \geq 3$ is Enriquez's outstanding
conjecture.  A positive resolution would extend Cycle 6 to
the full $\Mbar_{g,n}$ tower.

\textbf{(O3) Stokes classification at $g \geq 3$ confluent
strata.}
Jimbo-Miwa-Mori 1980 applies at rank $1$; the $\geq 3$-rank
higher-genus confluence requires Sabbah's irregular
Riemann-Hilbert correspondence, whose explicit form for the
$\Delta_{ij} \cap \partial_{\mathrm{non-sep}}\Mbar_{g,n}$
stratum is unknown.

\textbf{(O4) $W$-algebra extension of (G2).}
$\Pi_2$ of \texttt{chiral\_climax\_platonic.tex}: the
universal $W$-arena $\ConnConf^{W}$ replaces Arnold's
$\eta_{ij}$ with universal $W$-KZ forms.  The sl$_N$
principal case is accessible via Feigin-Frenkel screening;
non-simply-laced Langlands-dual side is conjectural.

\textbf{(O5) Universal BRST resolution functor.}
$\Pi_3$: the canonical BRST resolutions for standard
landscape families are family-specific; no universal
functor $\BRST \colon \ChirAlg^{\mathrm{Koszul}} \to
\ChirAlg^{\mathrm{free}}$ is known.  This is the G3-side
complement of G2's BRST identity $K = -c_{\mathrm{ghost}}$.

%% ====================================================================
\section{Hidden structure discovered during the attack}
%% ====================================================================
\label{sec:G2-hidden}

\textbf{(H-hidden-1) FM degeneration is a contracting homotopy,
not just a projector.}
Cycle 1 produced the FM contracting homotopy $h_{\mathrm{FM}}$
satisfying $[\dbar, h_{\mathrm{FM}}] = \mathrm{id} - p_{\mathrm{cogen}}$.
This is a stronger statement than the programme has previously
committed to: the bar complex retracts onto its cogenerators
via an explicit geometric contraction (the FM degeneration).
This gives the FM compactification a chain-level homotopical
content beyond its use as a boundary-extraction tool.
Cross-reference: \texttt{feedback\_homotopy\_retract\_is\_data.md}
(AP142).

\textbf{(H-hidden-2) Arnold three-term as deconcatenation identity.}
The reason strict chain-level equality holds on $T^c_n$
($n \geq 3$) is that the reorganisation from
\emph{adjacent-swap coderivation} (bar side) to
\emph{pair-sum Arnold} (KZ side) is governed by Arnold's
three-term relation.  This is not a mere combinatorial
identity; it is the homological content of the three-term
relation, reading it as a coderivation-vs-pair-sum
compatibility.  This suggests a deeper unification: the
Arnold three-term is the coherence condition for the
deconcatenation coproduct on $T^c$, restricted to
codim-two strata.

\textbf{(H-hidden-3) Stokes frontier at $g = 2$ separating
stratum is bounded.}
Cycle 4 found that the $g = 2$ separating stratum carries
\emph{at most rank-1} irregular Stokes structure; this is
lower than the naive expectation of rank depending on $n$.
The reason: the separating boundary is locally a product
of two $g = 1$ curves, on which the Bernard-Felder propagator
is regular, and the confluence $\Delta_{ij} \cap \partial_{\mathrm{sep}}$
localises to a $\mathbb{P}^1$ confluence of rank 1
(Jimbo-Miwa-Mori 1980).  This bounds the Vol II modular
bootstrap kernel at $g = 2$ in a useful way for the
$\Phi_{10}$-anchored BKM Class-A K3$\,\times\,$E identity
(UTI-2 at $d = 3$).

\textbf{(H-hidden-4) GRT$_1(\mathbb{Q})$ identifies Drinfeld's
six faces with seven.}
In Cycle 6, the GRT$_1$ torsor at $g = 0$ acts on 6 presentations
(F1-F6 in the classical Drinfeld sense).  The seventh face F7
(Wick / free-field $r$-matrix from CLAUDE.md essential constants)
is the image of a particular orbit element, not a new torsor
generator.  This refines the Vol~I $\S$2.3 statement.

%% ====================================================================
\section{Cycle log (raw)}
%% ====================================================================
\label{sec:G2-log}

\begin{enumerate}
\item \textsc{Cycle 1.}
Attack: chain-level vs homotopy equality.
Heal: strict chain-level on cogenerators (Step 1),
extended by Arnold three-term on $T^c_n$ (Step 2).
Witness: FM contracting homotopy $h_{\mathrm{FM}}$
satisfying $[\dbar, h_{\mathrm{FM}}] = \mathrm{id} - p_{\mathrm{cogen}}$.
\item \textsc{Cycle 2.}
Attack: strict vs functorial pullback.
Heal: KZ functoriality $\Rightarrow$ strict commutation;
cofree-coalgebra naturality $\Rightarrow$ morphism-intertwining.
\item \textsc{Cycle 3.}
Attack: globalisation pathologies across $g$.
Heal: three regimes: $g = 0$ unconditional,
$g = 1, 2$ unconditional for standard landscape at non-critical
level (Bernard-Felder, CEE), $g \geq 3$ F5 frontier.
\item \textsc{Cycle 4.}
Attack: rigor of Jimbo-Miwa at higher genus.
Heal: rank-1 rigorous via formal completion at
$\Delta_{ij} \cap \partial_{\mathrm{sep}}\Mbar_{2,n}$; higher
rank at $g \geq 3$ uses Sabbah 2008 (conjectural).
\item \textsc{Cycle 5.}
Attack: composition of chain homotopies between four presentations.
Heal: no separate homotopies needed -- all four are
projections of common FM residue datum; hexagon reduces to
identity arrows.
\item \textsc{Cycle 6.}
Attack: GRT$_1$ action at $g \geq 1$ -- theorem or conjecture?
Heal: theorem at $g \leq 2$ (Drinfeld 1990, Furusho 2010,
CEE 2009); conjecture at $g \geq 3$.
Arithmetic harmonic: Brown 2012 identifies GRT$_1(\mathbb{Q})$
with motivic Galois of MT($\mathbb{Z}$); $p = 691, 3617$
shadow-tower duality gives $p$-adic specialisations.
\end{enumerate}

\bigskip

\textbf{Remark on inscription.}
The healed Theorem~\ref{thm:G2-healed} of $\S$3 is ready for
inscription as a sharpened replacement of the existing scope
qualifiers in \texttt{chiral\_climax\_platonic.tex}
(Definition \ref{def:connconf-platonic}, around line 276-278)
and \texttt{genus\_2\_ddybe\_platonic.tex} (Cycle 4 Stokes
structure, around Theorem \ref{thm:fay-trisecant-genus-2-specific}).
The inscription lifts Conjecture
\ref{conj:climax-kzb-genus-one-platonic} ($\Pi_1$) to
\textbf{Theorem} at $g = 1, 2$ with the citations
Bernard 1988 / Felder 1998 / Calaque-Enriquez-Etingof 2009,
leaving $g \geq 3$ as the genuine open frontier.
The chain-level strict equality of Cycle 1 sharpens
Definition~\ref{def:pullback-identity-structural-platonic}
from ``modulo logarithmic gauge'' to ``strictly, with
coderivation-vs-pair-sum equivalence by Arnold three-term''.

\end{document}
